\chapter*{Предисловие}
\addcontentsline{toc}{chapter}{Предисловие}

Эта книга выросла из убеждения, что две великие традиции обучения с подкреплением\,---\,классическая математическая теория и современная практика выравнивания больших языковых моделей\,---\,заслуживают быть изложенными как единая, связная история. В последние годы идеи, десятилетиями существовавшие лишь на страницах теории управления и динамического программирования, вырвались в центр разработки систем искусственного интеллекта, лежа в основе обучения систем, которые пишут, рассуждают и ведут диалог. Именно этому пути посвящена настоящая книга.

\medskip

\textbf{Замечание об источниках вдохновения и оригинальности.}\quad
Каждый учебник опирается на труды предшественников, и эта книга не исключение. Мы хотим отдать должное двум книгам, которые сформировали облик всей области и, неизбежно, нашу собственную точку зрения,~--- одновременно разъяснив, чем настоящая книга от них отличается.

Первая~--- \emph{Reinforcement Learning: An Introduction} Ричарда С. Саттона и Эндрю Г. Барто, стандартный справочник по классическому RL\@. Саттон и Барто заложили педагогическую последовательность~--- от бандитов через МДП к методам временных разностей~--- которой с тех пор следует практически каждый курс. Наше изложение в частях~I и~II строится по той же логической дуге, но отличается акцентами и подачей: мы приводим более строгий анализ сходимости с учётом недавних результатов, включаем обобщённое оценивание преимущества (GAE) и PPO в основное повествование, а не в качестве дополнений, и мотивируем каждую классическую идею с прицелом на её будущую роль в обучении языковых моделей. Эти решения отражают уроки, извлечённые при разработке и преподавании университетских курсов по RL, а также при применении этих алгоритмов в производственных условиях.

Вторая~--- \emph{The RLHF Book} Натана Ламберта, наиболее полное из написанных практических руководств по обучению с подкреплением на основе обратной связи от людей для языковых моделей. Ламберт показал, что алгоритмические идеи, лежащие в основе ChatGPT, Claude и их преемников, поддаются точному и доступному изложению. Части~III и~IV нашей книги охватывают пересекающуюся тематику~--- моделирование вознаграждения, оптимизацию предпочтений, PPO для LLM~--- но расширяют обсуждение в направлениях, обусловленных нашим собственным исследовательским и инженерным опытом: мы даём единое формальное описание, связывающее DPO, IPO и KTO через RL-цель с KL-ограничением, и рассматриваем ORPO как самостоятельный подход без эталонной модели; посвящаем отдельные главы GRPO и обучению с подкреплением на верифицируемых вознаграждениях (RLVR)~--- темам, появившимся после первого издания книги Ламберта; а также включаем обширную главу о многошаговых агентных системах, которая перебрасывает мост между теорией RL и нарождающейся практикой языковых агентов, использующих инструменты.

Более широко: эта книга~--- не синтез двух упомянутых работ, а попытка рассказать \emph{новую} историю: единое непрерывное повествование, которое начинается с принципа оптимальности Беллмана и заканчивается обучающими пайплайнами, лежащими в основе современных моделей-рассуждения (reasoning models). Эта книга опирается на практический опыт масштабного обучения с RLHF, отладки моделей вознаграждения и практических трудностей, которые не может передать ни один сугубо теоретический текст. Там, где нам удаётся сделать этот мост ощущающимся естественным,~--- заслуга принадлежит сообществу, чьи открытые исследования сделали это возможным. Там, где мы не дотягиваем,~--- ответственность полностью лежит на нас.


\medskip

\textbf{Цели этой книги.}\quad
Центральная задача настоящего тома~--- выстроить единый, непрерывный мост между математическими основами, заложенными Беллманом, Саттоном и Барто, и практическими методами, которые сегодня определяют обучение и выравнивание больших языковых моделей. Путь от принципа оптимальности Беллмана до работающего RLHF-пайплайна короче, чем может показаться, но требует подлинного владения обоими мирами. Практик, знающий только инженерные рецепты, упускает теоретические гарантии и диагностическую интуицию, которую даёт теория. Теоретик, не соприкоснувшийся с реальностью обучения языковых моделей\,---\,нестабильностью моделей вознаграждения, чувствительностью гиперпараметров PPO, сложностью формализации предпочтений людей\,---\,плохо подготовлен к тому, чтобы вносить вклад в наиболее важное приложение RL на сегодняшний день.

Части~I и~II данной книги (охватывающие бандитов, МДП, динамическое программирование, методы Монте-Карло, обучение на основе временных разностей, аппроксимацию функций, градиенты стратегии и методы «актор-критик») написаны под влиянием строгого изложения Саттона и Барто и могут быть прочитаны полностью самостоятельно как законченный учебник по классическому RL. Части~III и~IV опираются на этот фундамент и охватывают моделирование вознаграждения, RLHF, прямую оптимизацию предпочтений (DPO), групповую оптимизацию относительной стратегии (GRPO), мультиагентные системы и открытые направления исследований, существенно используя практические идеи Ламберта.

\medskip

\textbf{Для кого эта книга.}\quad
Мы писали наш труд, имея в виду три пересекающиеся аудитории.

\emph{Аспиранты в области машинного обучения и ИИ}, желающие получить математически честное изложение обучения с подкреплением, которое не останавливается на классических результатах, а доводит их до границы выравнивания языковых моделей.

\emph{Инженеры и специалисты-практики в области МО}, знакомые с RLHF в контексте языковых моделей и желающие понять теоретический механизм, лежащий в основе реализации, чтобы уверенно отлаживать, расширять и улучшать его.

\emph{Исследователи, переходящие от классического обучения с учителем или без учителя к обучению с подкреплением}, включая тех, чьи основные интересы лежат в области выравнивания, безопасности или оценки,~--- и кому нужна прочная база в базовых концепциях перед погружением в исследовательскую литературу.

Книга не рассчитана на полных новичков в машинном обучении: предполагается некоторое предварительное знакомство с идеями этой области. Однако относительно обучения с подкреплением книга является самодостаточной: никаких предварительных знаний в области RL не требуется.

\medskip

\textbf{Необходимые знания.}\quad
Предполагается, что читатель уверенно владеет базовой теорией вероятностей (случайные величины, математическое ожидание, условная вероятность, теорема Байеса), линейной алгеброй (векторы, матрицы, собственные значения) и одно- и многомерным математическим анализом (производные, градиенты, правило цепочки). Знакомство с ключевыми концепциями обучения с учителем\,---\,функциями потерь, градиентным спуском, нейронными сетями, переобучением\,---\,сделает части~III и~IV значительно более доступными, хотя мы напоминаем ключевые идеи по мере необходимости. Знание Python полезно для реализационно-ориентированного материала в приложениях, но не является обязательным для понимания теоретического изложения.

\medskip

\textbf{Как читать эту книгу.}\quad
Четыре части книги образуют естественную последовательность, однако читатели с разной подготовкой могут начинать с разных мест.

\emph{Часть~I (Основы)} развивает классическую теорию с нуля: многорукие бандиты, марковские процессы принятия решений, динамическое программирование, методы Монте-Карло и обучение на основе временных разностей. Читателю без предварительного знакомства с RL следует прорабатывать эти главы по порядку.

\emph{Часть~II (Аппроксимация функций и глубокое RL)} переходит к задачам, где табличные методы уже неприменимы. Она охватывает линейную и нелинейную аппроксимацию функции ценности, методы градиента стратегии и семейство «актор-критик», завершая изложение алгоритмом PPO. Вместе с частью~I этот материал представляет собой законченное введение в классическое и глубокое RL.

\emph{Часть~III (RL для языковых моделей, главы~10--14)} применяет всё изложенное в частях~I и~II к современной задаче обучения и выравнивания больших языковых моделей. Читатели, уже владеющие классическим RL, могут переходить непосредственно сюда. В этих главах рассматриваются формирование вознаграждения и модели вознаграждения (глава~10), пайплайн обучения языковых моделей от предобучения до SFT (глава~11), прямая оптимизация предпочтений и её варианты (глава~12), групповая оптимизация относительной стратегии и онлайн-RL с верифицируемыми вознаграждениями (глава~13), а также инженерная реализация RLHF в промышленном масштабе (глава~14).

\emph{Часть~IV (Горизонты, главы~15--17)} исследует передний край области: рассуждение, использование инструментов и дообучение с подкреплением (глава~15); мультиагентные системы (глава~16); обзор открытых проблем и новых направлений (глава~17).

Каждая глава завершается набором упражнений, диапазон которых охватывает как простые проверочные задания, так и более сложные доказательства и задачи проектирования. Мы рекомендуем выполнять хотя бы несколько упражнений к каждой главе: обучение с подкреплением~--- это, что вполне символично, предмет, который лучше всего понимается через практику.

\medskip

\textbf{Сопроводительный код и ноутбуки.}\quad
К книге прилагается открытый репозиторий на GitHub, содержащий практические Jupyter-ноутбуки~--- по одному на каждую главу. В этих ноутбуках представлены работающие реализации всех основных алгоритмов, обсуждаемых в тексте,~--- от $\varepsilon$-жадных бандитов и Q-обучения до DQN, PPO, DPO и GRPO,~--- а также визуализации и управляемые эксперименты. Все ноутбуки рассчитаны на работу на бесплатных платформах, таких как Google Colab или Kaggle, и не требуют специализированного оборудования. Репозиторий также включает подробные решения упражнений с полными доказательствами и выводами, а также инфраструктуру для переводов на другие языки силами сообщества.

Сопроводительный репозиторий доступен по адресу:

\medskip
\centerline{\url{https://github.com/pyshka501/rl-textbook}}
\medskip

\noindent
Мы рекомендуем работать с ноутбуками параллельно с изучением соответствующих глав. Сочетание математической строгости текста и практического экспериментирования с кодом является, по нашему опыту, наиболее эффективным способом выработать подлинное мастерство в области обучения с подкреплением.

\medskip

Мы надеемся, что эта книга послужит вам хорошим проводником~--- где бы вы ни находились на пути от классической теории к современной практике выравнивания.

\bigskip

\hfill\textit{Москва, 2026}
